% Method section (draft). Numbers: A1 / A1b ablation (paper/tables/ablation.tex), A2 action sweep
% (results/experiments/action_sweep/summary.md), A3 stall and trade-off (results/stats/a3_stall.csv,
% results/figures/a3_*). Settings: scripts/experiment.py POLICIES["final"].
\section{The proposed method: \ours{}}
\label{sec:method}

\ours{} keeps the evolutionary framework of IMOEA-DQN~\cite{jin2026imoeadqn} wherever the DAG model allows and changes
it in three places. First, the convergence ranking of the three-tier environmental selection, Jin et al.'s interval
dominance degree, is replaced by a decomposition rank that does not lose its selection pressure once most of the
population is mutually non-dominated (Sections~\ref{sec:stall}--\ref{sec:rankc}). Second, parents mate with
neighbours in objective space (Section~\ref{sec:variation}). Third, the selection weights $(\alpha, \beta, \gamma)$,
which IMOEA-DQN chooses every generation with a deep Q-network, are fixed to a value chosen on training instances
(Section~\ref{sec:weights}); Section~\ref{sec:learned} shows that the learned controller does not beat it.
Algorithm~\ref{alg:main} gives the main loop.

\subsection{Encoding and variation}
\label{sec:variation}
A solution is the pair $(\mathbf{x}, \mathbf{o})$ of Section~\ref{sec:model}. The initial population draws every
allocation gene uniformly from the nodes the task's device can reach and every priority string from a randomised
Kahn's algorithm, which picks uniformly among the ready tasks. Algorithm~\ref{alg:var} creates one child per member of
the population. The partner of member $i$ is drawn from its $T$ nearest neighbours by Euclidean distance between
normalised interval midpoints, the mating restriction of MOEA/D~\cite{zhang2007moead}. Crossover keeps the
dual-adaptive scheme of Jin et al.: the crossover probability is 0.7 for the better half of the population by
uncertainty front and 0.9 for the other half, and the allocation string is recombined gene by gene (DAEC), taking each
gene from either parent with equal probability. Jin et al.'s simulated binary crossover cannot be applied to a task
order, so the priority strings are recombined by precedence-preserving order-based crossover
(POX)~\cite{bierwirth1996permutation}: a random binary mask decides, position by position, which parent supplies its
leftmost task not yet used, and any mask yields a topological order. Mutation reassigns each allocation gene with
probability $1/K$ to another reachable node and moves each task with probability $1/K$ to a random position after its
last predecessor and before its first successor. Every child is therefore feasible by construction.

\subsection{Three-tier environmental selection}
Jin et al.'s selection ranks the merged population $P \cup Q$ of $2N$ solutions in three ways:
\begin{itemize}
\item \emph{Rank C} (convergence): the interval dominance degree $R_{IN}(x) = \sum_{y \prec_{IN} x} S(y)$, where $S(y)$ is
  the number of pool members $y$ interval-dominates; lower is better.
\item \emph{Rank D} (diversity): the interval crowding distance $\mathrm{ICD}(x) = \sum_{y} \mathrm{ICD}(x, y)$ over the
  $k = \mathrm{round}(\sqrt{2N})$ most similar solutions $y$, with
  $\mathrm{ICD}(x, y) = \sum_m |m(\tilde F_m(x)) - m(\tilde F_m(y))| / (\mu(x, y) + V(x) + V(y) + 1)$, where $V$ is the
  volume of a solution's objective box and $\mu$ the overlap volume of two boxes; higher is better.
\item \emph{Rank U} (uncertainty): the dominance degree $R_U$ under Pareto dominance on $(U_1, U_2)$; lower is better.
\end{itemize}
Given $(\alpha, \beta, \gamma)$ with $\alpha + \beta + \gamma = 1$, the next population takes the best $\mathrm{round}(\alpha N)$
solutions by Rank C, then the best $\mathrm{round}(\beta N)$ not yet chosen by Rank D, and the rest by Rank U
(Algorithm~\ref{alg:sel}). Exact copies of a chosen genotype are skipped, so clones cannot fill the population.

\subsection{Why the interval dominance degree stalls}
\label{sec:stall}
$R_{IN}$ separates solutions only while some of them dominate others. Under interval dominance two solutions are
incomparable whenever their intervals overlap enough in some objective, so after a few hundred generations most of the
merged population has $R_{IN} = 0$. When more than $\alpha N$ pool members share the best value, Rank C no longer
selects: the $\alpha N$ convergence places are filled by the tie-break, which Jin et al.\ leave unspecified and which
is random in a literal implementation. Figure~\ref{fig:a3_stall} measures this on two instances. With $R_{IN}$, the
share of Rank C places decided by the tie-break rises to 86--100\% after generation 100 and to almost 100\% after
generation 500, when 117--121 of the 210 pool members share $R_{IN} = 0$ for $\alpha N = 88$ places. The convergence tier
then draws at random from the first front, and the population drifts along it instead of approaching the Pareto front.
The effect grows with the size of the DAG: the fixed-action Jin backbone has a mean IGD of 0.052 on the 100-task
workflows of Suite~B and 0.175 on the 1000-task ones (Section~\ref{sec:results}).

\subsection{Tchebycheff Rank C}
\label{sec:rankc}
\ours{} replaces $R_{IN}$ by a rank that orders every solution. Let $W$ be a set of $N$ Das--Dennis weight vectors and let
$G^L, G^U$ be the objectives scaled per objective to $[0, 1]$ over the pool. For each $\mathbf{w} \in W$ the interval
Tchebycheff value of a solution is
\begin{equation}
\tilde g(x \mid \mathbf{w}) = \Big[\max_m w_m G^L_m(x), \; \max_m w_m G^U_m(x)\Big],
\label{eq:tch}
\end{equation}
and the pool is sorted by it. Sorting intervals by the possibility degree~\eqref{eq:possibility} at threshold 0.5 is
sorting by midpoint, which is how IMOEA/D compares two solutions on a subproblem. Let $\pi_{\mathbf{w}}(x) \in
\{0, \dots, 2N - 1\}$ be the position of $x$ when the pool is sorted by ascending $m(\tilde g(\cdot \mid \mathbf{w}))$.
The new Rank C of a solution is its best position over all these orders,
\begin{equation}
R_C(x) = \min_{\mathbf{w} \in W} \pi_{\mathbf{w}}(x),
\label{eq:rankc}
\end{equation}
so the best solution of every subproblem has $R_C = 0$, and a solution is ranked well when it is close to the best on at
least one direction. This is the ranking of MSOPS~\cite{hughes2003msops} applied to interval midpoints. What \ours{}
adds is its use as the convergence tier inside the dominance-, diversity- and uncertainty-aware selection of Jin et al.,
with ties broken by descending interval crowding distance rather than at random. With this rank the share of
convergence places decided by the tie-break stays at 11--23\% throughout the run (Fig.~\ref{fig:a3_stall}). Unlike
MOEA/D-style methods, \ours{} keeps a generational population and Jin et al.'s diversity and uncertainty tiers.

\subsection{Selection weights}
\label{sec:weights}
The weights $(\alpha, \beta, \gamma)$ are fixed to $(5/6, 1/6, 0)$, which gives 88 places by Rank C and 17 by Rank D in
a population of 105. The value was chosen before any test run by holding each of the 28 actions of IMOEA-DQN's action
set fixed on 36 training instances that are separate from the test instances (three further instances of each Suite~A
configuration, with data $\times 1$ and $\times 20$). It had the best mean IGD rank (2.25 of 28); on average its IGD was 1.13 times
the best action's on each instance, and 1.71 times at worst. Only four actions, all with $\alpha \ge 2/3$, were ever
best on an instance. The other weights trade convergence for precision along one curve
(Fig.~\ref{fig:a3_tradeoff}): a small uncertainty share, $(5/6, 0, 1/6)$, lowers the imprecision per solution by 10\%
for an IGD ratio of 1.15 instead of 1.13, while more weight on Rank D or Rank U raises IGD steeply. The neighbourhood
size $T = 10$ and the probability of mating with a neighbour, 1.0, were tuned on the same training instances.

% Computational complexity (R2). Derived from the code (gat_imoea_dqn/operators.py, selection.py,
% intervals.py, evaluator.py); measured times in Table~\ref{tab:timing} (scripts/build_paper.py).
\subsection{Computational complexity}
\label{sec:complexity}

Let $N$ be the population size, $M = 3$ the number of objectives, $K = |V|$ and $|E|$ the numbers of tasks
and data dependencies, and $F$ the number of interval non-dominated fronts in a population
($F \le N$, and small in practice because most of the population soon shares the first front).
One generation of \ours{} costs the following.

\begin{itemize}
\item \emph{Variation} (Algorithm~\ref{alg:var}). The interval dominance matrix of $P$ takes $O(MN^2)$ and peeling
  its fronts $O(FN^2)$; the uncertainty fronts on $(U_1, U_2)$ take the same. The neighbourhoods need all pairwise
  midpoint distances, $O(MN^2)$, and a sort of each row, $O(N^2 \log N)$. Building one child takes $O(K + |E|)$:
  DAEC and POX are linear in $K$, a shift mutation moves a task within its precedence window in $O(K)$, and with
  mutation rate $1/K$ one task moves per child in expectation. In total $O\big(N^2(M + F + \log N) + N(K + |E|)\big)$.
\item \emph{Evaluation} (Algorithm~\ref{alg:eval}). The interval critical-path calculus visits every task and every
  dependency once, $O(K + |E|)$ per solution and $O\big(N(K + |E|)\big)$ per generation.
\item \emph{Selection} (Algorithm~\ref{alg:sel}) on the pool of $2N$. The Tchebycheff values for $N$ weight vectors take
  $O(MN^2)$ and sorting the pool once per vector $O(N^2 \log N)$; the interval crowding distance needs all pairs,
  $O(MN^2)$; the uncertainty dominance degree $O(N^2)$; filling the three tiers with genotype de-duplication $O(NK)$.
  In total $O\big(N^2(M + \log N) + NK\big)$.
\end{itemize}

A generation therefore costs $O\big(N^2(M + F + \log N) + N(K + |E|)\big)$, and a run of $G_{\max}$ generations
$O\big(G_{\max}[N^2(M + F + \log N) + N(K + |E|)]\big)$ time with $O(N^2 + NK)$ memory. Jin et al.'s Rank C, the
interval dominance degree $R_{IN}$, costs $O(MN^2)$ for the dominance matrix plus $O(N^2)$ for the degrees, so the
Tchebycheff Rank C adds only the $O(N^2 \log N)$ of the per-vector sorts; neighbourhood mating adds the same order.
Neither changes the number of evaluations, which is $N(G_{\max} + 1)$ for \ours{} and every baseline.

With $N = 105$ the population terms are about $N^2(M + \log_2 2N) \approx 1.2 \times 10^5$ elementary operations per
generation, while the evaluation term $N(K + |E|)$ grows from about $3 \times 10^4$ for the 100-task Cybershake
workflow ($|E| = 192$) to $3 \times 10^5$ at 1000 tasks and $6 \times 10^5$ at 2000. The measured times
(Table~\ref{tab:timing}) show that \ours{} needs about as long as IMOEA/D and the fixed-action Jin backbone on every
instance, and that its time grows only 1.8-fold from 100 to 2000 Cybershake tasks: the evaluation loop is compiled and
has a small constant, so the population terms, which do not depend on $K$, account for most of the run time. The
learned controller of IMOEA-DQN adds a state computation and a network forward pass every generation and a training
step every 50 generations; on the same instances this costs 2--52\% more time than the fixed-action backbone, without
improving the result (see the analysis of learned control in the paper).

